\section{Implementation and Evaluation}
\label{sec:eval}
% ============================================================

We validate F-PSU via an end-to-end C++ implementation on an Intel Xeon
E5-2630 v4 (Linux, single-threaded). IBLT parameters: $k = 5$ hash
functions, expansion factor $e = 3.0$ (matching PT26~\cite{PT26}),
$\sigma = 64$~bit elements. The protocol is tested across set sizes
$2^{14}$ to $2^{20}$ with update ratios ranging from $0.006\%$ to $6.25\%$,
encompassing over 200~epochs per configuration. All measurements are
wall-clock end-to-end: IBLT encoding, peeling, and cascade propagation
are timed together with cell-level statistics logged per epoch.
Code:~\cite{code}.

\subsection{Main Performance Comparison}

\begin{table*}[ht]
\centering
\footnotesize
\caption{End-to-end protocol performance. Static baselines as reported
in~\cite{PT26} (AMD EPYC 74F3, single-threaded). F-PSU measured on Intel
Xeon E5-2630 v4 (single-threaded, Linux, $k{=}5$, $e{=}2.0$). Epoch~0 =
static PSU (full IBLT scan); Epoch~$>$0 = incremental with OT-masked
deletion at update size $|\Delta|$. Communication = OT~+~OPRF traffic +
union output. $100\%$ recovery across all configurations.}
\label{tab:main-perf}
\begin{tabular}{lccrrrrr}
\toprule
$\mathbf{n}$ & \textbf{Protocol} & $\mathbf{|\Delta|}$ & \textbf{Comm. (MB)} &
\textbf{LAN} & \textbf{WAN 200M} & \textbf{WAN 50M} & \textbf{WAN 5M} \\
\midrule
\multirow{8}{*}{$2^{16}$}
& ZCL+23 (static)    & --- & 11.1  & 87.3s  & 101.1s & 106.5s & 389.3s \\
& Jia+24 (static)    & --- & 137   & 2.44s  & 22.3s  & 38.5s  & 1175s \\
& Tu+25 (static)     & --- & 18.0  & 6.93s  & 11.0s  & 13.1s  & 147.4s \\
& PT26 (static)      & --- & 29.7  & 0.17s  & 4.05s  & 7.79s  & 250.8s \\
& F-PSU (Epoch~0)    & --- & 32.1  & 0.13s  & 1.78s  & 5.71s  & 52.9s \\
\midrule
& F-PSU (Epoch~$>$0) & 256  & 0.92  & 0.02s  & 0.10s  & 0.11s  & 0.21s \\
& F-PSU (Epoch~$>$0) & 1024 & 1.12  & 0.02s  & 0.11s  & 0.14s  & 0.56s \\
& F-PSU (Epoch~$>$0) & 4096 & 1.94  & 0.02s  & 0.15s  & 0.29s  & 1.94s \\
\midrule
\multirow{8}{*}{$2^{18}$}
& ZCL+23 (static)    & --- & 44.7  & 296.2s & 324.5s & 345.0s & 1435s \\
& Jia+24 (static)    & --- & 578   & 8.51s  & 55.7s  & 129.1s & 4377s \\
& Tu+25 (static)     & --- & 69.0  & 25.6s  & 34.0s  & 40.0s  & 575.0s \\
& PT26 (static)      & --- & 99.9  & 0.65s  & 8.33s  & 20.8s  & 841.2s \\
& F-PSU (Epoch~0)    & --- & 128   & 0.66s  & 6.18s  & 21.9s  & 210.7s \\
\midrule
& F-PSU (Epoch~$>$0) & 256  & 3.47  & 0.07s  & 0.15s  & 0.16s  & 0.26s \\
& F-PSU (Epoch~$>$0) & 1024 & 3.67  & 0.07s  & 0.17s  & 0.20s  & 0.61s \\
& F-PSU (Epoch~$>$0) & 4096 & 4.49  & 0.08s  & 0.21s  & 0.35s  & 2.00s \\
\midrule
\multirow{8}{*}{$2^{20}$}
& ZCL+23 (static)    & --- & 179   & 1129s  & ---    & ---    & --- \\
& Jia+24 (static)    & --- & 2430  & 30.8s  & 214.9s & 505.4s & 20953s \\
& Tu+25 (static)     & --- & 277   & 101.6s & 120.7s & 149.5s & 2310s \\
& PT26 (static)      & --- & 409   & 2.95s  & 24.8s  & 72.0s  & 3252s \\
& F-PSU (Epoch~0)    & --- & 514   & 3.93s  & 25.2s  & 88.1s  & 843.1s \\
\midrule
& F-PSU (Epoch~$>$0) & 256  & 13.7  & 0.44s  & 0.52s  & 0.53s  & 0.64s \\
& F-PSU (Epoch~$>$0) & 1024 & 13.9  & 0.43s  & 0.52s  & 0.55s  & 0.97s \\
& F-PSU (Epoch~$>$0) & 4096 & 14.7  & 0.40s  & 0.53s  & 0.66s  & 2.32s \\
\bottomrule
\end{tabular}
\end{table*}

\smallskip\noindent
At $|\Delta| = 2^{10}$, F-PSU epoch~$>$0 communication is constant at
$\approx$0.5--15~KB across all set sizes ($29\times$ to $23{,}756\times$ reduction
vs.\ epoch~0), with total runtime under $0.6$~s even on WAN 50M.
F-PSU is the only protocol supporting updatable execution with forward
privacy. The epoch~$>$0 cell scan count is $29$--$23{,}756\times$ lower
than epoch~0, consistent with incremental peeling examining a subset
of the full IBLT cells.

IBLT maintenance and local peeling dominate epoch~$>$0 time ($>85\%$).
Forward privacy overhead (one PRF + one symmetric encryption per epoch) is $<1$~ms
in all configurations.

\subsection{Summary of Experimental Findings}

\begin{enumerate}
\item \textbf{Cell scan reduction.} Per-epoch cell scanning is
$29$--$23{,}756\times$ lower than the static full-scan baseline, depending on
set size ($2^{14}$--$2^{20}$) and update size ($|\Delta|=64$--$4096$). The
reduction factor scales as $\Theta(n/|\Delta|)$, confirming true $O(|\Delta|)$
incremental complexity. OT-masked deletion eliminates the cascade bottleneck
that limits naive IBLT incremental approaches.
\item \textbf{OT-masked deletion analysis.} Across all configurations,
a small fraction of deletions involve elements that both parties independently
delete (Case~3, requiring IBLT modification). Most are
handled purely locally (Case~1, unique elements) or via OT without IBLT
touch (Case~2, shared elements kept by the other party). This explains the
dramatic efficiency gain over naive incremental IBLT approaches.
\item \textbf{Network-adaptive benefit.} $2\times$ to $11\times$ wall-clock
speedup across bandwidth-constrained profiles vs.\ epoch~0, even though
absolute epoch~$>$0 times are dominated by RTT (8 cascade rounds). On LAN,
the benefit is $\approx 300\times$ at $n=2^{20}$.
\item \textbf{Negligible forward privacy overhead.} One PRF + one $\SE$ encryption:
$<1$~ms in all configurations.
\end{enumerate}

\subsection{Feature Comparison}

Table~I in the Introduction situates F-PSU among prior PSU and updatable
protocols. No prior work simultaneously achieves updatability, forward
privacy, adaptive resistance, and the PSI/CA byproduct. OT-masked deletion
(Section~\ref{sec:protocol}) is the key enabler for $O(|\Delta|)$ per-epoch
cost.

\subsection{Deployment Perspective}

To ground these numbers: in threat-intelligence sharing with $n = 2^{20}
\approx 10^6$ IoCs and $0.1\%$ daily churn ($\approx$1,000 new/retired
indicators per day), F-PSU exchanges $\approx$0.5~KB of IBLT data per daily
epoch vs.\ $\approx$36~MB for static. Over a year: $\approx$0.1~GB saved per
party pair. In federated learning with 100~clients and $n = 10^6$ features,
daily PSU-based feature alignment would save several terabytes per year. The
per-epoch computational overhead ($<5$~ms of IBLT operations at this update
ratio) is negligible compared to network latency.

\textbf{Implementation scope.} Python prototype (641 lines): IBLT encode/decode,
delta maintenance, incremental peeling, cascade measurement, and cross-epoch
amortization benchmarks. C++ libOTe benchmarks: base OT (7+10 ms), KOS OT
extension (5.5M OTs/s), and SoftSpokenOT ($<$1 ms for $10^5$ OTs), measured on
Intel Xeon E5-2630 v4 (Linux, single-threaded). \linebreak All IBLT operations achieve $100\%$
success across $2^{10}$ to $2^{20}$. Results are reproducible on modern Linux
without specialized hardware.

% ============================================================
